\section{Related Work}
\label{sec:related}

\paragraph{Active inference and the free-energy principle.}
Friston's free-energy principle frames perception, learning, and
action as the minimization of a variational free energy that
decomposes into an accuracy term (how well internal states predict
observations) and a complexity term (how far the posterior drifts
from the prior)~\citep{friston2010}. Its computational instantiation
as variational message passing~\citep{parr2017} and gradient-flow
inference~\citep{daCosta2020} provides the theoretical scaffolding
for our functional. Crucially, classical active-inference
implementations use Gaussian or factorized categorical posteriors,
both of which admit closed-form KL divergences. Our choice of a
\emph{product of simplices} $\rho \in (\Delta^{32})^4$ preserves the
closed-form complexity term per species but forces numerical
handling of the accuracy term, which we discuss in
\S\ref{sec:query-conditioned-f}.

\paragraph{JEPA and self-prediction.}
Joint Embedding Predictive Architecture~\citep{lecun2022} learns
representations by predicting target embeddings from context
embeddings, with subsequent realizations on images
\citep{assran2023ijepa} and video~\citep{bardes2024vjepa}. Our use of
$g_{\mathrm{JEPA}}$ as the accuracy decoder of an active-inference
free energy interprets the predictor as an amortized posterior
$q(s \mid o)$ over psycholinguistic simplex latents. In contrast to
standard JEPA, we (i) fix the target domain (the text embedding
manifold) rather than learning it jointly, and (ii) force the
context ($\rho$) onto a structured simplex. This restriction is what
makes the free-energy framing meaningful: the KL complexity term
would be ill-defined on an unconstrained embedding.

\paragraph{Wasserstein gradient flows and JKO.}
The Jordan-Kinderlehrer-Otto scheme~\citep{jko1998} discretizes
Wasserstein gradient flows on probability measures, providing the
dynamic substrate for our oracle. Recent neural-network
instantiations~\citep{bunne2022wgf} train surrogates that mimic JKO
trajectories, but typically without text conditioning. Adjacent work
on text-to-distribution bridging \citep{rezende2015variational} uses
affine flows parameterized by text, but on free-support densities,
not on bounded simplex latents. To the best of our knowledge,
\textsc{QueryConditionedF} is the first JKO formulation that
(a) couples to a text encoder via a JEPA-style accuracy decoder,
(b) operates on a product-of-simplices latent, and (c) preserves the
full Friston accuracy/complexity decomposition.

\paragraph{Text encoders on structured latents.}
Most work on text-to-latent mapping targets an unconstrained
embedding space (e.g., sentence-BERT families~\citep{reimers2019sbert})
or discrete code spaces via VQ-VAE~\citep{vandenoord2017neural}.
Two notable exceptions are softmax-normalized topic-model bridges
(LDA-style distributions over a simplex) and categorical latents in
dialogue modeling, but neither frames the mapping as an active-inference
step. Our work sits in the intersection: a text encoder whose output
drives a dynamical process on a simplex, where the dynamics are
justified by the free-energy principle rather than by likelihood
maximization.

\paragraph{French psycholinguistic resources.}
Our heuristic labeler relies on Lexique.org~\citep{newboucard2004lexique}
for word frequency bins, SpaCy-fr~\citep{honnibal2020spacy} for
dependency parsing, and the \texttt{phonemizer}
toolkit~\citep{phonemizer} for IPA extraction.
Levelt's working-memory architecture~\citep{levelt1989speaking}
motivates the four-species decomposition (phono, sem, lex,
syntax), which Baddeley and Hitch's multi-component working memory
model \citep{baddeley1974} reinforces by treating phonological and
semantic stores as architecturally separable. We inherit the
four-way split from the codebase but emphasize that the
\textsc{QueryConditionedF} machinery is not tied to this particular
decomposition --- any finite species count with pre-defined priors
would fit the same functional.

\paragraph{Substrate-agnostic bridges.}
Our broader research programme treats cross-modal coupling as a
substrate-agnostic communication problem, following the Nerve
Protocol~\citep{saillant2026nervewml} which decouples protocol from
substrate via discrete codebook plus phase-amplitude
multiplexing. \textsc{QueryConditionedF} can be viewed as a
specialization: the text encoder is the ``source substrate'' and the
psycholinguistic simplex is the ``target substrate,'' with the JEPA
decoder playing the role of a translator. This framing is made
explicit in our companion full paper (in preparation).
